\documentclass[aspectratio=169]{beamer}
\usetheme{Madrid}
\usecolortheme{whale}
\usepackage[utf8]{inputenc}
\usepackage[T1]{fontenc}
\usepackage[spanish]{babel}
\usepackage{tikz}
\usetikzlibrary{shapes.geometric, arrows.meta, positioning, fit, backgrounds}
\usepackage{booktabs}
\usepackage{fontawesome5}
\usepackage{xcolor}
\usepackage{array}
\usepackage{colortbl}

% ── Paleta de colores ──────────────────────────────────────────────
\definecolor{gcpBlue}{HTML}{1A73E8}
\definecolor{gcpGreen}{HTML}{34A853}
\definecolor{gcpYellow}{HTML}{FBBC04}
\definecolor{gcpRed}{HTML}{EA4335}
\definecolor{localGreen}{HTML}{2E7D32}
\definecolor{localBg}{HTML}{E8F5E9}
\definecolor{gcpBg}{HTML}{E3F2FD}
\definecolor{darkBg}{HTML}{1E2A38}

\setbeamercolor{palette primary}{bg=gcpBlue,fg=white}
\setbeamercolor{palette secondary}{bg=gcpBlue!80,fg=white}
\setbeamercolor{palette tertiary}{bg=gcpBlue!60,fg=white}
\setbeamercolor{frametitle}{bg=gcpBlue,fg=white}
\setbeamercolor{title}{fg=white}
\setbeamercolor{structure}{fg=gcpBlue}
\setbeamerfont{frametitle}{size=\large,series=\bfseries}

% ── Estilos TikZ ──────────────────────────────────────────────────
\tikzset{
  box/.style={
    rectangle, rounded corners=4pt,
    minimum width=3.2cm, minimum height=0.6cm,
    text centered, font=\scriptsize\bfseries,
    draw=#1!80, fill=#1!15, text=#1!90
  },
  boxW/.style={box=gcpBlue, minimum width=2.8cm},
  boxL/.style={box=localGreen},
  boxG/.style={box=gcpBlue},
  boxN/.style={box=darkBg, fill=gray!10},
  decision/.style={
    diamond, aspect=2.2,
    minimum width=2.8cm, minimum height=0.7cm,
    text centered, font=\tiny\bfseries,
    draw=gcpYellow!80!black, fill=gcpYellow!20
  },
  arrow/.style={-Stealth, thick, color=gray!70},
  arrowB/.style={-Stealth, thick, color=gcpBlue},
  arrowG/.style={-Stealth, thick, color=localGreen},
  annot/.style={font=\tiny, text=gray!70, inner sep=1pt}
}

% ── Metadatos ─────────────────────────────────────────────────────
\title[Sistema de Facturas PDF/XML]{\textbf{Sistema de Automatización de Facturas}\\\large Arquitectura técnica y flujo operativo}
\subtitle{Alupar | La Virgen · Proyecto 2}
\author{Equipo de Automatización}
\date{Mayo 2026}

\begin{document}

% ─────────────────────────────────────────────────────────────────
% PORTADA
% ─────────────────────────────────────────────────────────────────
\begin{frame}
\titlepage
\end{frame}

% ─────────────────────────────────────────────────────────────────
% AGENDA
% ─────────────────────────────────────────────────────────────────
\begin{frame}{Contenido}
\tableofcontents
\end{frame}

\section{Flujo de procesamiento}
\section{Recursos GCP}
\section{Rol de Workato}
\section{Costos estimados}

% ─────────────────────────────────────────────────────────────────
% SLIDE: FLUJO COMPLETO (parte 1)
% ─────────────────────────────────────────────────────────────────
\begin{frame}{Flujo de procesamiento · Entrada y validación}
\begin{center}
\begin{tikzpicture}[node distance=0.42cm and 1.1cm,
  scale=0.82, every node/.append style={transform shape}]

  % Origen
  \node[boxW] (workato) {Workato / Postman};

  % Intake
  \node[boxN, below=of workato] (intake) {POST /api/intake};

  % Auth
  \node[decision, below=of intake] (auth) {¿Secret válido?};

  % 401
  \node[box=gcpRed, right=1.6cm of auth] (r401) {401 Unauthorized};

  % Tipo
  \node[decision, below=of auth] (ftype) {Tipo de archivo};

  % Rejected
  \node[box=gcpRed, right=1.6cm of ftype] (rej) {rejected++};

  % Hash
  \node[boxN, below=of ftype] (hash) {SHA-256 del buffer};

  % Dedup
  \node[decision, below=of hash] (dedup) {¿Ya existe?};

  % Skip
  \node[box=gcpGreen, right=1.6cm of dedup] (skip) {Devuelve existente};

  % Guardar archivo
  \node[boxN, below=of dedup] (save) {Guardar archivo original};

  % Flechas
  \draw[arrow] (workato) -- (intake);
  \draw[arrow] (intake)  -- (auth);
  \draw[arrow] (auth)    -- node[annot, left]{sí} (ftype);
  \draw[arrow] (auth)    -- node[annot, above]{no} (r401);
  \draw[arrow] (ftype)   -- node[annot, left]{pdf/xml} (hash);
  \draw[arrow] (ftype)   -- node[annot, above]{unknown} (rej);
  \draw[arrow] (hash)    -- (dedup);
  \draw[arrow] (dedup)   -- node[annot, left]{no} (save);
  \draw[arrow] (dedup)   -- node[annot, above]{sí} (skip);

  % Leyenda
  \begin{scope}[on background layer]
    \node[fill=localBg, rounded corners, opacity=0.5,
      fit=(workato)(intake)(auth)(ftype)(hash)(dedup)(save)(r401)(rej)(skip),
      inner sep=4pt] {};
  \end{scope}

\end{tikzpicture}
\end{center}
\end{frame}

% ─────────────────────────────────────────────────────────────────
% SLIDE: FLUJO COMPLETO (parte 2 - bifurcación Local vs GCP)
% ─────────────────────────────────────────────────────────────────
\begin{frame}{Flujo de procesamiento · Extracción y persistencia}
\begin{columns}[T]

% ── Columna LOCAL ──────────────────────────────────
\begin{column}{0.47\textwidth}
  \begin{block}{\textcolor{localGreen}{\faLaptopCode\ \ Entorno LOCAL}}
    \begin{center}
    \begin{tikzpicture}[node distance=0.52cm]
      \node[boxL] (ls)  {storage/raw/<id>/};
      \node[boxL, below=of ls] (lp)  {pdf-parse + regex};
      \node[boxL, below=of lp] (lx)  {xml2js (XML)};
      \node[boxL, below=of lx] (lc)  {classifyDocument};
      \node[boxL, below=of lc] (li)  {inferConcept};
      \node[boxL, below=of li] (ld)  {data/documents.json};
      \draw[arrowG] (ls) -- (lp);
      \draw[arrowG] (lp) -- (lx);
      \draw[arrowG] (lx) -- (lc);
      \draw[arrowG] (lc) -- (li);
      \draw[arrowG] (li) -- (ld);
    \end{tikzpicture}
    \end{center}
  \end{block}
\end{column}

% ── Columna GCP ────────────────────────────────────
\begin{column}{0.47\textwidth}
  \begin{block}{\textcolor{gcpBlue}{\faCloud\ \ Entorno GCP (Cloud Run)}}
    \begin{center}
    \begin{tikzpicture}[node distance=0.52cm]
      \node[boxG] (gs)  {Cloud Storage (RAW bucket)};
      \node[boxG, below=of gs] (gp)  {Document AI · Invoice Processor};
      \node[boxG, below=of gp] (gx)  {xml2js (XML)};
      \node[boxG, below=of gx] (gc)  {classifyDocument};
      \node[boxG, below=of gc] (gi)  {inferConcept};
      \node[boxG, below=of gi] (gd)  {Firestore · colección documents};
      \draw[arrowB] (gs) -- (gp);
      \draw[arrowB] (gp) -- (gx);
      \draw[arrowB] (gx) -- (gc);
      \draw[arrowB] (gc) -- (gi);
      \draw[arrowB] (gi) -- (gd);
    \end{tikzpicture}
    \end{center}
  \end{block}
\end{column}

\end{columns}

\vspace{0.3cm}
\centering
\small\textbf{Respuesta 202 →} \texttt{\{requestId, accepted, rejected, documents[\,]\}}
\end{frame}

% ─────────────────────────────────────────────────────────────────
% SLIDE: FLUJO POST-PROCESAMIENTO
% ─────────────────────────────────────────────────────────────────
\begin{frame}{Flujo de procesamiento · Salidas operativas}
\begin{center}
\begin{tikzpicture}[node distance=0.6cm and 2cm]

  \node[boxN] (fs)  {Documento procesado en Firestore};
  \node[boxN, below=of fs] (dash) {Dashboard web (Cloud Run)};

  \node[boxG, below left=0.8cm and 1.8cm of dash, align=center, text width=3.2cm]  (tab) {Pestañas por tipo\\ Factura · Comprobante · Nota};
  \node[boxG, below=0.8cm of dash, align=center, text width=2.6cm]                 (det) {Vista de detalle\\ + PDF original};
  \node[boxG, below right=0.8cm and 1.8cm of dash, align=center, text width=3.4cm] (zip) {POST /api/exports/:concepto\\ $\rightarrow$ ZIP en Cloud Storage};

  \draw[arrowB] (fs)   -- (dash);
  \draw[arrowB] (dash) -- (tab);
  \draw[arrowB] (dash) -- (det);
  \draw[arrowB] (dash) -- (zip);

\end{tikzpicture}
\end{center}

\vspace{0.3cm}
\begin{columns}
  \begin{column}{0.45\textwidth}
    \begin{exampleblock}{Estados de documento}
      \texttt{pendiente} → \texttt{procesado} → \texttt{error}
    \end{exampleblock}
  \end{column}
  \begin{column}{0.45\textwidth}
    \begin{exampleblock}{Dedupe automático}
      Hash SHA-256 del contenido del archivo evita duplicados en cada corrida
    \end{exampleblock}
  \end{column}
\end{columns}
\end{frame}

% ─────────────────────────────────────────────────────────────────
% SLIDE: RECURSOS GCP
% ─────────────────────────────────────────────────────────────────
\begin{frame}{Recursos GCP utilizados}
\begin{columns}[T]
\begin{column}{0.55\textwidth}
  \rowcolors{2}{gcpBg}{white}
  \arrayrulecolor{black}
  \begin{tabular}{@{}p{3.2cm}p{5cm}@{}}
    \toprule
    \textbf{Servicio} & \textbf{Rol en el sistema} \\
    \midrule
    \textbf{Cloud Run} & Servidor principal. Corre la API Node.js y el dashboard web. Escala a 0 cuando no hay tráfico. \\
    \textbf{Cloud Build} & Compila y empaqueta la imagen Docker en cada deploy. \\
    \textbf{Container Registry} & Almacena la imagen Docker del servicio. \\
    \textbf{Cloud Storage} & Guarda PDFs/XMLs originales (bucket raw) y ZIPs exportados (bucket exports). \\
    \textbf{Firestore} & Base de datos de documentos procesados. Alimenta el dashboard en tiempo real. \\
    \textbf{Document AI} & Extrae campos estructurados de PDFs de facturas (número, fecha, monto, emisor). \\
    \textbf{IAM} & Service account con roles mínimos para Storage, Firestore y Document AI. \\
    \bottomrule
  \end{tabular}
\end{column}
\begin{column}{0.4\textwidth}
  \begin{alertblock}{Arquitectura serverless}
    Cloud Run escala a cero instancias cuando no hay tráfico. Solo se cobra por requests procesados.
  \end{alertblock}

  \vspace{0.4cm}

  \begin{block}{Buckets configurados}
    \footnotesize
    \texttt{invoice-automation-497420-raw}\\
    \quad → PDFs/XMLs originales\\[4pt]
    \texttt{invoice-automation-497420-exports}\\
    \quad → ZIPs por concepto
  \end{block}

  \vspace{0.4cm}

  \begin{block}{Document AI · Processor}
    \footnotesize
    Tipo: \texttt{INVOICE\_PROCESSOR}\\
    Ubicación: \texttt{us}\\
    ID: \texttt{5624412ac0f58ebf}
  \end{block}
\end{column}
\end{columns}
\end{frame}

% ─────────────────────────────────────────────────────────────────
% SLIDE: ROL DE WORKATO
% ─────────────────────────────────────────────────────────────────
\begin{frame}{Rol de Workato en el flujo}
\begin{columns}[T]
\begin{column}{0.5\textwidth}
  \textbf{Lo que hace Workato:}
  \vspace{0.2cm}
  \begin{itemize}
    \item[\faEnvelope] Monitorea el buzón de correo (Outlook / Gmail) en forma continua.
    \item[\faPaperclip] Detecta correos con adjuntos PDF o XML.
    \item[\faFilter] Filtra por remitente, tamaño, extensión y lista de permitidos.
    \item[\faDownload] Descarga el adjunto en binario (bytes).
    \item[\faPaperPlane] Envía \texttt{multipart/form-data} al endpoint\\
      \texttt{POST /api/intake} de Cloud Run con:
      \begin{itemize}
        \item[\scriptsize\faFile] Archivo(s) en campo \texttt{files}
        \item[\scriptsize\faTag] Metadatos: \texttt{messageId}, \texttt{sender}, \texttt{subject}, \texttt{receivedAt}
        \item[\scriptsize\faKey] Header \texttt{x-workato-secret}
      \end{itemize}
    \item[\faSync] Gestiona reintentos automáticos ante errores de red.
    \item[\faBell] Envía alertas ante fallos reiterados.
  \end{itemize}
\end{column}
\begin{column}{0.45\textwidth}
  \textbf{Lo que Workato \underline{no} hace:}
  \vspace{0.2cm}
  \begin{itemize}
    \item[\faTimes] Leer ni interpretar el contenido del PDF/XML.
    \item[\faTimes] Clasificar tipo de documento.
    \item[\faTimes] Extraer campos (número, monto, fecha).
    \item[\faTimes] Guardar datos en base de datos.
    \item[\faTimes] Generar ZIPs ni reportes.
  \end{itemize}

  \vspace{0.5cm}

  \begin{exampleblock}{Contrato entre Workato y GCP}
    \footnotesize
    \textbf{Workato envía:} archivo + metadatos del correo\\[3pt]
    \textbf{GCP responde:} \texttt{202 Accepted} con campos extraídos, tipo documental y concepto\\[3pt]
    \textbf{Idempotencia:} el hash del archivo evita reprocesar duplicados aunque Workato reintente
  \end{exampleblock}
\end{column}
\end{columns}
\end{frame}

% ─────────────────────────────────────────────────────────────────
% SLIDE: COSTOS ESTIMADOS GCP
% ─────────────────────────────────────────────────────────────────
\begin{frame}{Costos estimados GCP}
\begin{columns}[T]
\begin{column}{0.58\textwidth}
  \rowcolors{2}{gcpBg}{white}
  \arrayrulecolor{black}
  \begin{tabular}{@{}p{2.8cm}p{2.2cm}p{3cm}@{}}
    \toprule
    \textbf{Servicio} & \textbf{USD/mes} & \textbf{Base del cálculo} \\
    \midrule
    Cloud Run       & \$0 -- \$15  & \textasciitilde500 req/día · 512 MB RAM \\
    Cloud Build     & \$0 -- \$5   & 120 min/día gratis · builds ocasionales \\
    Container Registry & \$1 -- \$3 & imagen \textasciitilde200 MB \\
    Cloud Storage   & \$2 -- \$10  & 5--50 GB PDFs/XMLs/ZIPs \\
    Firestore       & \$1 -- \$15  & 10k escrituras + 50k lecturas/día \\
    Document AI     & \$15 -- \$60 & 500--2000 páginas/mes · \$0.03/pág \\
    \midrule
    \textbf{Total estimado} & \textbf{\$20 -- \$108} & Operación normal \\
    \bottomrule
  \end{tabular}

  \vspace{0.3cm}
  \footnotesize
  \textit{Precios referenciales · región us-central1 · mayo 2026.}\\
  \textit{Document AI: primeras 1000 pág/mes son gratis (Free Tier).}
\end{column}

\begin{column}{0.38\textwidth}
  \begin{alertblock}{\faExclamationTriangle\ \ Variable clave}
    El costo de \textbf{Document AI} domina el total y depende directamente del volumen de páginas procesadas por mes.
  \end{alertblock}

  \vspace{0.4cm}

  \begin{block}{Escenarios}
    \footnotesize
    \begin{tabular}{ll}
      \textbf{Bajo}  & \textless 500 pág/mes \\
                     & \textasciitilde \$20/mes \\[4pt]
      \textbf{Medio} & 500--2000 pág/mes \\
                     & \textasciitilde \$50/mes \\[4pt]
      \textbf{Alto}  & \textgreater 2000 pág/mes \\
                     & \textasciitilde \$100+/mes \\
    \end{tabular}
  \end{block}

  \vspace{0.3cm}

  \begin{exampleblock}{Free Tier aplicable}
    \footnotesize
    Cloud Run: 2M req/mes gratis\\
    Firestore: 1 GB almacenamiento\\
    Document AI: 1000 pág/mes
  \end{exampleblock}
\end{column}
\end{columns}
\end{frame}

% ─────────────────────────────────────────────────────────────────
% SLIDE: RESUMEN
% ─────────────────────────────────────────────────────────────────
\begin{frame}{Resumen de la solución}
\begin{columns}[T]
\begin{column}{0.48\textwidth}
  \begin{block}{Capa de captura · Workato}
    \begin{itemize}
      \item Monitor de correo Outlook/Gmail
      \item Filtrado y descarga de adjuntos PDF/XML
      \item Envío a Cloud Run con metadatos
      \item Reintentos y alertas operativas
    \end{itemize}
  \end{block}

  \vspace{0.3cm}

  \begin{block}{Capa de procesamiento · GCP}
    \begin{itemize}
      \item Cloud Run: API + Dashboard web
      \item Document AI: extracción de campos
      \item Firestore: persistencia y consultas
      \item Cloud Storage: archivos y exports
    \end{itemize}
  \end{block}
\end{column}

\begin{column}{0.48\textwidth}
  \begin{exampleblock}{Capacidades habilitadas}
    \begin{itemize}
      \item[\faCheckCircle] Clasificación automática de documentos
      \item[\faCheckCircle] Extracción de campos clave (número, fecha, monto, emisor)
      \item[\faCheckCircle] Dashboard por pestañas (Facturas · Comprobantes · Notas)
      \item[\faCheckCircle] Exportación ZIP por concepto
      \item[\faCheckCircle] Dedupe por hash de contenido
      \item[\faCheckCircle] Trazabilidad de estado por documento
      \item[\faCheckCircle] Soporte PDF y XML en el mismo flujo
    \end{itemize}
  \end{exampleblock}

  \vspace{0.3cm}

  \begin{alertblock}{Entorno de prueba disponible}
    \tiny\ttfamily https://proyecto2-facturas-\\lanrfivtaq-uc.a.run.app
  \end{alertblock}
\end{column}
\end{columns}
\end{frame}

\end{document}
